\documentclass[tikz,border=10pt]{standalone}
\usepackage{tikz}
\usepackage{amsmath}
\usepackage{amssymb}
\usepackage{ctex}
\usepackage{pgfplots}
\pgfplotsset{compat=1.18}
\usetikzlibrary{calc, positioning, arrows.meta, decorations.pathreplacing}

\begin{document}
\begin{tikzpicture}[>=Stealth]

%% ============================
%%  左子图：实异根 — p(λ)=(λ-2)(λ-5)，根=2,5
%% ============================
\begin{scope}[xshift=0cm]

  \begin{axis}[
    name=ax1,
    width=5.8cm, height=5.0cm,
    xmin=-0.5, xmax=7.0,
    ymin=-5.5, ymax=8.0,
    axis lines=center,
    xlabel={$\lambda$},
    ylabel={$p(\lambda)$},
    xlabel style={font=\small, right},
    ylabel style={font=\small, above},
    xtick={2,5},
    xticklabels={$2$,$5$},
    ytick={0},
    tick label style={font=\scriptsize},
    clip=false,
  ]
    % 曲线 p(λ)=(λ-2)(λ-5)
    \addplot[blue!70!black, thick, domain=-0.3:6.8, samples=80]
      {(x-2)*(x-5)};

    % 零点标注（谱点）
    \addplot[red!70!black, only marks, mark=*, mark size=2pt]
      coordinates {(2,0) (5,0)};

    % 谱标注
    \node[font=\scriptsize, red!70!black, above] at (axis cs:2,0.3) {$\lambda=2$};
    \node[font=\scriptsize, red!70!black, above] at (axis cs:5,0.3) {$\lambda=5$};

    % 谱集合标签
    \node[font=\scriptsize, fill=red!8, draw=red!50, rounded corners=2pt,
          inner sep=3pt, align=center]
      at (axis cs:3.5,6.5) {谱 $\sigma(A)=\{2,5\}$};

  \end{axis}

  % 子标题
  \node[font=\small\bfseries, align=center] at ([yshift=-0.7cm]ax1.south)
    {实异根\\$p(\lambda)=(\lambda-2)(\lambda-5)$};

\end{scope}

%% ============================
%%  中子图：重根 — p(λ)=(λ-3)²，切线点λ=3
%% ============================
\begin{scope}[xshift=7.2cm]

  \begin{axis}[
    name=ax2,
    width=5.8cm, height=5.0cm,
    xmin=0.0, xmax=6.0,
    ymin=-1.0, ymax=7.0,
    axis lines=center,
    xlabel={$\lambda$},
    ylabel={$p(\lambda)$},
    xlabel style={font=\small, right},
    ylabel style={font=\small, above},
    xtick={3},
    xticklabels={$3$},
    ytick={0},
    tick label style={font=\scriptsize},
    clip=false,
  ]
    % 曲线 p(λ)=(λ-3)²
    \addplot[orange!80!black, thick, domain=0.2:5.8, samples=80]
      {(x-3)^2};

    % 切线点（重根）
    \addplot[red!70!black, only marks, mark=*, mark size=2pt]
      coordinates {(3,0)};
    \node[font=\scriptsize, red!70!black, above right] at (axis cs:3,0.15)
      {重根 $\lambda=3$};

    % 切线 y=0（λ轴相切）
    \addplot[gray!50, dashed, domain=1.0:5.0] {0};

    % 谱集合标签
    \node[font=\scriptsize, fill=orange!10, draw=orange!50, rounded corners=2pt,
          inner sep=3pt, align=center]
      at (axis cs:2.5,5.5) {谱 $\sigma(A)=\{3\}$\\（代数重数=2）};

  \end{axis}

  % 子标题
  \node[font=\small\bfseries, align=center] at ([yshift=-0.7cm]ax2.south)
    {重根（切点）\\$p(\lambda)=(\lambda-3)^2$};

\end{scope}

%% ============================
%%  右子图：复根 — p(λ)=λ²-4λ+5，无实根，谱={2±i}
%%  下方配小型复平面示意
%% ============================
\begin{scope}[xshift=14.4cm]

  \begin{axis}[
    name=ax3,
    width=5.8cm, height=5.0cm,
    xmin=-0.5, xmax=5.5,
    ymin=-0.5, ymax=5.5,
    axis lines=center,
    xlabel={$\lambda$},
    ylabel={$p(\lambda)$},
    xlabel style={font=\small, right},
    ylabel style={font=\small, above},
    xtick={2},
    xticklabels={$2$},
    ytick={1},
    yticklabels={$1$},
    tick label style={font=\scriptsize},
    clip=false,
  ]
    % 曲线 p(λ)=λ²-4λ+5=(λ-2)²+1，最小值1，无实根
    \addplot[purple!70!black, thick, domain=-0.2:5.2, samples=80]
      {x^2 - 4*x + 5};

    % 最小值点（顶点）
    \addplot[gray!60, only marks, mark=o, mark size=2pt]
      coordinates {(2,1)};
    \node[font=\scriptsize, gray!60, right] at (axis cs:2.1,1.2)
      {顶点 $(2,1)$};

    % 无实根提示
    \node[font=\scriptsize, purple!80, fill=purple!8, draw=purple!40,
          rounded corners=2pt, inner sep=3pt, align=center]
      at (axis cs:3.8,4.0) {无实根\\曲线不过 $x$ 轴};

  \end{axis}

  % 复平面示意（轴外下方）
  \begin{scope}[shift={([xshift=0.4cm, yshift=-2.8cm]ax3.south west)}]
    \draw[->,gray!60,thin] (-0.3,0) -- (3.5,0)
      node[right,font=\scriptsize,gray] {$\mathrm{Re}$};
    \draw[->,gray!60,thin] (1.4,-0.8) -- (1.4,1.5)
      node[above,font=\scriptsize,gray] {$\mathrm{Im}$};

    % 复数特征值 2+i, 2-i
    \fill[purple!70!black] (2.0, 0.7) circle (2pt)
      node[right,font=\scriptsize,purple!70!black] {$2+i$};
    \fill[purple!70!black] (2.0,-0.7) circle (2pt)
      node[right,font=\scriptsize,purple!70!black] {$2-i$};

    % 虚线连接到实轴
    \draw[dashed,purple!30] (2.0,0.7) -- (2.0,-0.7);
    \draw[dashed,purple!20] (2.0,0.7) -- (2.0,0);
    \fill[gray!50] (2.0,0) circle (1.5pt)
      node[below,font=\scriptsize,gray] {$2$};

    \node[font=\scriptsize\bfseries,purple!70!black] at (0.5,1.2)
      {复平面};
    \node[font=\scriptsize,fill=purple!8,draw=purple!40,rounded corners=2pt,
          inner sep=3pt,align=center]
      at (0.5,0.1) {谱 $\sigma(A)$\\$=\{2\pm i\}$};
  \end{scope}

  % 子标题
  \node[font=\small\bfseries, align=center] at ([yshift=-0.7cm]ax3.south)
    {复根（无实根）\\$p(\lambda)=\lambda^2-4\lambda+5$};

\end{scope}

%% 整体标题
\node[font=\large\bfseries] at (9.8, 5.4)
  {特征多项式的根与谱};

\end{tikzpicture}
\end{document}
